\clearpage
\begin{abstract}
Does lowering numeric precision reduce the GPU energy of large language model (LLM) inference, and can routing prompts by complexity exploit any such saving? We measured this with the NVML energy counter of NVIDIA RTX 6000 Ada GPUs on a shared academic cluster, across two inference stacks, two model sizes and two batch sizes.

The two stacks gave different answers. With bitsandbytes in Hugging Face Transformers and a LoRA adapter attached to every tier, Llama-3.2-1B at batch size 1 used about as much energy per token at 4-bit NF4 as at float16 (within about 7\% in each run) and about 1.8 to 1.9 times as much at 8-bit. Under vLLM, FP8 and AWQ 4-bit cut energy per token for Llama-3.1-8B-Instruct by 38\% and 64\% at both batch sizes, and FP8 cut it by 36-37\% for Llama-3.2-1B. Because the stacks differ in engine, kernel and quantization scheme, and 4-bit was tested on different models, these results describe serving configurations and do not isolate bit-width. The accuracy cost is unsettled: tinyGSM8k fell by 7.2 points (FP8) and 14.1 points (AWQ) under flexible scoring and by 6.6 and 5.7 points under strict scoring, but a paired check on 75 GSM8K prompts, using a coarser measure, found no comparable fall.

Repeated float16 measurements ranged from 1.59 to 12.97 J/token, a factor of about 8. Three standard safeguards accepted the inflated values, whereas throughput and mean power exposed them, and we propose a continuous monitoring protocol. Green-Weight, a five-feature complexity sensor with a Mamdani fuzzy controller, performed no better than a random router with the same tier mix at any tested threshold, and its features predicted sensitivity to quantization at about chance level. On the 8B precision pool, static AWQ cost 240 J/request against 477 J/request for the router.

\textbf{Keywords:} LLM inference, energy measurement, quantization, FP8, AWQ, vLLM, bitsandbytes, measurement validity, fuzzy routing.

\end{abstract}
